\documentclass[10pt,a4paper]{amsart}
\usepackage{amsmath,amssymb,amsthm,graphicx}
% Retained minimal preamble for the Stacks Project chapter more-morphisms.tex (More on Morphisms).
% Verbatim excerpt of the chapter preamble: Xy-pic setup (native preamble lines 14-18) and the native
% theorem declarations with the mathematical macros (native preamble lines 154-217). Copied mechanically;
% nothing added, removed or redefined. The wrapper class and the amsmath/amssymb/amsthm/graphicx base
% packages are supplied by the assembler, so the class/externaldocument part of the native preamble is omitted.
% For commutative diagrams we use Xy-pic
\usepackage[all]{xy}

% We use 2cell for 2-commutative diagrams.
\xyoption{2cell}
% Theorem environments.
%
\theoremstyle{plain}
\newtheorem{theorem}[subsection]{Theorem}
\newtheorem{proposition}[subsection]{Proposition}
\newtheorem{lemma}[subsection]{Lemma}

\theoremstyle{definition}
\newtheorem{definition}[subsection]{Definition}
\newtheorem{example}[subsection]{Example}
\newtheorem{exercise}[subsection]{Exercise}
\newtheorem{situation}[subsection]{Situation}

\theoremstyle{remark}
\newtheorem{remark}[subsection]{Remark}
\newtheorem{remarks}[subsection]{Remarks}

\numberwithin{equation}{subsection}

% Macros
%
\def\lim{\mathop{\mathrm{lim}}\nolimits}
\def\colim{\mathop{\mathrm{colim}}\nolimits}
\def\Spec{\mathop{\mathrm{Spec}}}
\def\Hom{\mathop{\mathrm{Hom}}\nolimits}
\def\Ext{\mathop{\mathrm{Ext}}\nolimits}
\def\SheafHom{\mathop{\mathcal{H}\!\mathit{om}}\nolimits}
\def\SheafExt{\mathop{\mathcal{E}\!\mathit{xt}}\nolimits}
\def\Sch{\mathit{Sch}}
\def\Mor{\mathop{\mathrm{Mor}}\nolimits}
\def\Ob{\mathop{\mathrm{Ob}}\nolimits}
\def\Sh{\mathop{\mathit{Sh}}\nolimits}
\def\NL{\mathop{N\!L}\nolimits}
\def\CH{\mathop{\mathrm{CH}}\nolimits}
\def\proetale{{pro\text{-}\acute{e}tale}}
\def\etale{{\acute{e}tale}}
\def\QCoh{\mathit{QCoh}}
\def\Ker{\mathop{\mathrm{Ker}}}
\def\Im{\mathop{\mathrm{Im}}}
\def\Coker{\mathop{\mathrm{Coker}}}
\def\Coim{\mathop{\mathrm{Coim}}}

% Boxtimes
%
\DeclareMathSymbol{\boxtimes}{\mathbin}{AMSa}{"02}

%
% Macros for moduli stacks/spaces
%
\def\QCohstack{\mathcal{QC}\!\mathit{oh}}
\def\Cohstack{\mathcal{C}\!\mathit{oh}}
\def\Spacesstack{\mathcal{S}\!\mathit{paces}}
\def\Quotfunctor{\mathrm{Quot}}
\def\Hilbfunctor{\mathrm{Hilb}}
\def\Curvesstack{\mathcal{C}\!\mathit{urves}}
\def\Polarizedstack{\mathcal{P}\!\mathit{olarized}}
\def\Complexesstack{\mathcal{C}\!\mathit{omplexes}}
% \Pic is the operator that assigns to X its picard group, usage \Pic(X)
% \Picardstack_{X/B} denotes the Picard stack of X over B
% \Picardfunctor_{X/B} denotes the Picard functor of X over B
\def\Pic{\mathop{\mathrm{Pic}}\nolimits}
\def\Picardstack{\mathcal{P}\!\mathit{ic}}
\def\Picardfunctor{\mathrm{Pic}}
\def\Deformationcategory{\mathcal{D}\!\mathit{ef}}
\usepackage{cleveref}
\crefname{originalreference}{original reference}{original references}
\Crefname{originalreference}{Original reference}{Original references}
\makeatletter
\newlabel{descent-lemma-flat-finitely-presented-permanence}{{02KL}{source}{}{sourcepointer-1}{}}
\newlabel{descent-lemma-flat-finitely-presented-permanence@cref}{{[originalreference][0][]02KL}{[1][0][]source}{}{}{}}
\makeatother
\crefname{originallemma}{original lemma}{original lemmas}
\Crefname{originallemma}{Original lemma}{Original lemmas}
\makeatletter
\newlabel{lemma-composition-pseudo-coherent}{{0681}{source}{}{sourcepointer-2}{}}
\newlabel{lemma-composition-pseudo-coherent@cref}{{[originallemma][0][]0681}{[1][0][]source}{}{}{}}
\makeatother
\crefname{originallemma}{original lemma}{original lemmas}
\Crefname{originallemma}{Original lemma}{Original lemmas}
\makeatletter
\newlabel{lemma-dominate-fppf}{{05WN}{source}{}{sourcepointer-3}{}}
\newlabel{lemma-dominate-fppf@cref}{{[originallemma][0][]05WN}{[1][0][]source}{}{}{}}
\makeatother
\crefname{originallemma}{original lemma}{original lemmas}
\Crefname{originallemma}{Original lemma}{Original lemmas}
\makeatletter
\newlabel{lemma-flat-finite-presentation-pseudo-coherent}{{0695}{source}{}{sourcepointer-4}{}}
\newlabel{lemma-flat-finite-presentation-pseudo-coherent@cref}{{[originallemma][0][]0695}{[1][0][]source}{}{}{}}
\makeatother
\crefname{originallemma}{original lemma}{original lemmas}
\Crefname{originallemma}{Original lemma}{Original lemmas}
\makeatletter
\newlabel{lemma-permanence-pseudo-coherent}{{0683}{source}{}{sourcepointer-5}{}}
\newlabel{lemma-permanence-pseudo-coherent@cref}{{[originallemma][0][]0683}{[1][0][]source}{}{}{}}
\makeatother
\crefname{originalreference}{original reference}{original references}
\Crefname{originalreference}{Original reference}{Original references}
\makeatletter
\newlabel{more-algebra-lemma-finite-extension-pseudo-coherent}{{0673}{source}{}{sourcepointer-6}{}}
\newlabel{more-algebra-lemma-finite-extension-pseudo-coherent@cref}{{[originalreference][0][]0673}{[1][0][]source}{}{}{}}
\makeatother
\crefname{originalreference}{original reference}{original references}
\Crefname{originalreference}{Original reference}{Original references}
\makeatletter
\newlabel{more-algebra-lemma-summands-relative-pseudo-coherent}{{0676}{source}{}{sourcepointer-7}{}}
\newlabel{more-algebra-lemma-summands-relative-pseudo-coherent@cref}{{[originalreference][0][]0676}{[1][0][]source}{}{}{}}
\makeatother
\title{Pseudo-coherence of a morphism is fppf local on the source}
\author{The Stacks Project Authors}
\date{stacks-project:more-morphisms.tex:local-snapshot-20261011; GNU FDL 1.2 or later}
\begin{document}
\maketitle
\noindent\textit{Editorial note.} Counted span is the author's own proof environment, source lines 17626-17677, immediately following the statement at 17620-17624.\par\medskip
\noindent\textit{Editorial note.} The author states the remaining algebra problem explicitly ('the algebra problem is the following: ...') inside the span; it is retained verbatim.\par\medskip
\noindent\textit{Editorial note.} Three cross-chapter pointers (Descent, More on Algebra) are recorded in assembly.external\_references.\par\medskip
\noindent\textit{Allowed standard prerequisites (editorial).}\begin{itemize}
\item Dominating an fppf covering by a nicer one. chapter lemma lemma-dominate-fppf, cited at proof line 17630. Replaces the given fppf covering by one with a Zariski covering of the source and affine pieces, which makes the algebra question accessible.
\item Composition and flat finite presentation inputs. chapter lemmas lemma-composition-pseudo-coherent and lemma-flat-finite-presentation-pseudo-coherent, cited in the proof. Controls pseudo-coherence under composition and under flat morphisms of finite presentation.
\item Permanence of pseudo-coherence. chapter lemma lemma-permanence-pseudo-coherent with Descent, lemma descent-lemma-flat-finitely-presented-permanence, cited in the proof. Moves pseudo-coherence along the flat finitely presented pieces of the covering.
\item Algebraic ingredients on pseudo-coherence of ring maps. More on Algebra, lemmas more-algebra-lemma-finite-extension-pseudo-coherent and more-algebra-lemma-summands-relative-pseudo-coherent, cited at proof lines 17669-17675. Solves the resulting ring-theoretic problem: if R -> A' is pseudo-coherent and A' is a free A-module of finite rank, then R -> A is pseudo-coherent.
\end{itemize}\medskip
\ifdefined\hypertarget\hypertarget{sourcepointer-1}{}\fi
\noindent\textit{Source reference.} Original reference 02KL: focused-seeds/geometry/S9-morph-pseudo-coherent-fppf-source/ref-tags.txt, source lines 4--4. Navigation pointer only: the counted proof cites this statement by \textbackslash{}ref as a standing input; the deduction the counted proof has to make is carried out inside the counted span, and no part of that deduction is delegated to this pointer.\par\medskip
\ifdefined\hypertarget\hypertarget{sourcepointer-2}{}\fi
\noindent\textit{Source reference.} Original lemma 0681: sources/stacks-more-morphisms/pinned-original-source.tex, source lines 17402--17402. Navigation pointer only: the counted proof cites this statement by \textbackslash{}ref as a standing input; the deduction the counted proof has to make is carried out inside the counted span, and no part of that deduction is delegated to this pointer.\par\medskip
\ifdefined\hypertarget\hypertarget{sourcepointer-3}{}\fi
\noindent\textit{Source reference.} Original lemma 05WN: sources/stacks-more-morphisms/pinned-original-source.tex, source lines 13649--13649. Navigation pointer only: the counted proof cites this statement by \textbackslash{}ref as a standing input; the deduction the counted proof has to make is carried out inside the counted span, and no part of that deduction is delegated to this pointer.\par\medskip
\ifdefined\hypertarget\hypertarget{sourcepointer-4}{}\fi
\noindent\textit{Source reference.} Original lemma 0695: sources/stacks-more-morphisms/pinned-original-source.tex, source lines 17428--17428. Navigation pointer only: the counted proof cites this statement by \textbackslash{}ref as a standing input; the deduction the counted proof has to make is carried out inside the counted span, and no part of that deduction is delegated to this pointer.\par\medskip
\ifdefined\hypertarget\hypertarget{sourcepointer-5}{}\fi
\noindent\textit{Source reference.} Original lemma 0683: sources/stacks-more-morphisms/pinned-original-source.tex, source lines 17440--17440. Navigation pointer only: the counted proof cites this statement by \textbackslash{}ref as a standing input; the deduction the counted proof has to make is carried out inside the counted span, and no part of that deduction is delegated to this pointer.\par\medskip
\ifdefined\hypertarget\hypertarget{sourcepointer-6}{}\fi
\noindent\textit{Source reference.} Original reference 0673: focused-seeds/geometry/S9-morph-pseudo-coherent-fppf-source/ref-tags.txt, source lines 5--5. Navigation pointer only: the counted proof cites this statement by \textbackslash{}ref as a standing input; the deduction the counted proof has to make is carried out inside the counted span, and no part of that deduction is delegated to this pointer.\par\medskip
\ifdefined\hypertarget\hypertarget{sourcepointer-7}{}\fi
\noindent\textit{Source reference.} Original reference 0676: focused-seeds/geometry/S9-morph-pseudo-coherent-fppf-source/ref-tags.txt, source lines 6--6. Navigation pointer only: the counted proof cites this statement by \textbackslash{}ref as a standing input; the deduction the counted proof has to make is carried out inside the counted span, and no part of that deduction is delegated to this pointer.\par\medskip
\begin{lemma}
\label{lemma-pseudo-coherent-fppf-local-source}
The property $\mathcal{P}(f) =$``$f$ is pseudo-coherent''
is fppf local on the source.
\end{lemma}

\begin{proof}
Let $f : X \to S$ be a morphism of schemes.
Let $\{g_i : X_i \to X\}$ be an fppf covering such that each composition
$f \circ g_i$ is pseudo-coherent. According to
Lemma \ref{lemma-dominate-fppf}
there exist
\begin{enumerate}
\item a Zariski open covering $X = \bigcup U_j$,
\item surjective finite locally free morphisms $W_j \to U_j$,
\item Zariski open coverings $W_j = \bigcup_k W_{j, k}$,
\item surjective finite locally free morphisms $T_{j, k} \to W_{j, k}$
\end{enumerate}
such that the fppf covering $\{h_{j, k} : T_{j, k} \to X\}$ refines the given
covering $\{X_i \to X\}$. Denote $\psi_{j, k} : T_{j, k} \to X_{\alpha(j, k)}$
the morphisms that witness the fact that $\{T_{j, k} \to X\}$ refines
the given covering $\{X_i \to X\}$. Note that $T_{j, k} \to X$ is a flat,
locally finitely presented morphism, so both $X_i$ and $T_{j, k}$ are
pseudo-coherent over $X$ by
Lemma \ref{lemma-flat-finite-presentation-pseudo-coherent}.
Hence $\psi_{j, k} : T_{j, k} \to X_i$ is pseudo-coherent, see
Lemma \ref{lemma-permanence-pseudo-coherent}.
Hence $T_{j, k} \to S$ is pseudo coherent as the composition
of $\psi_{j, k}$ and $f \circ g_{\alpha(j, k)}$, see
Lemma \ref{lemma-composition-pseudo-coherent}.
Thus we see we have reduced the lemma to the case of
a Zariski open covering (which is OK) and the case of a covering
given by a single surjective finite locally free morphism which we
deal with in the following paragraph.

\medskip\noindent
Assume that $X' \to X \to S$ is a sequence of morphisms of schemes
with $X' \to X$ surjective finite locally free and $X' \to Y$ pseudo-coherent.
Our goal is to show that $X \to S$ is pseudo-coherent.
Note that by
Descent, Lemma \ref{descent-lemma-flat-finitely-presented-permanence}
the morphism $X \to S$ is locally of finite presentation.
It is clear that the problem reduces to the case that $X'$, $X$ and $S$
are affine and $X' \to X$ is free of some rank $r > 0$. The corresponding
algebra problem is the following: Suppose $R \to A \to A'$ are ring maps
such that $R \to A'$ is pseudo-coherent, $R \to A$ is of finite presentation,
and $A' \cong A^{\oplus r}$ as an $A$-module. Goal: Show $R \to A$ is
pseudo-coherent. The assumption that $R \to A'$ is pseudo-coherent
means that $A'$ as an $A'$-module is pseudo-coherent relative to $R$. By
More on Algebra,
Lemma \ref{more-algebra-lemma-finite-extension-pseudo-coherent}
this implies that $A'$ as an $A$-module is pseudo-coherent relative to $R$.
Since $A' \cong A^{\oplus r}$ as an $A$-module we see that
$A$ as an $A$-module is pseudo-coherent relative to $R$, see
More on Algebra,
Lemma \ref{more-algebra-lemma-summands-relative-pseudo-coherent}.
This by definition means that $R \to A$ is pseudo-coherent and we win.
\end{proof}

\end{document}
